\section{为什么 inference 是一门完整的系统工程}

前面的 CS25 课程大量讨论了模型能力、训练规模与新架构；本讲把视角移动到模型真正被用户调用的地方。核心问题不是“怎样把一个 checkpoint 放进 API”这么窄，而是：\textbf{怎样把应用需求翻译成可测量的 workload 与 correctness contract，再沿着模型、引擎、GPU、网络、容器、监控和优化逐层实现？} 这条链上任一层定义错误，都可能让更快的 kernel 变成无效劳动。

\lecturefigure{slide-002.jpg}{从训练转向推理：模型最终要被人使用}{官方 deck 第 2 页；课堂 00:01:55--00:02:42}

\begin{importantbox}{本讲的 inference working definition}
推理服务不是一次前向传播，而是一个持续接收请求、管理 token 与 KV cache、调度 CPU/GPU 工作、满足延迟和可靠性目标、记录可调试证据并控制成本的在线或批处理系统。模型 forward 只是数据路径的一部分；SLO、排队、缓存、故障恢复与 eval 同样属于系统语义。
\end{importantbox}

\teachervoice{00:00:56--00:01:55，老师说训练研究回答“模型从哪里来”，而 inference 回答“模型最终怎样被人使用”。即使读者只想做训练，也应理解自己的架构、tokenizer、精度和并行选择会怎样约束后续 serving。}

\subsection{从 cost center 到可交付系统}

训练把资金、数据与算力转换成权重，但权重本身并不自动形成一个可收费、可依赖、可维护的产品。推理把权重包进交互接口、业务流程与可靠性机制；同时，现代 post-training 会大量调用当前策略生成 rollout，再把结果回流到优化过程，因此“训练”和“推理”也不再是完全分开的两套基础设施。

\lecturefigure{slide-003.jpg}{为什么 inference 会吸引更多工程资源}{官方 deck 第 3 页；课堂 00:02:42--00:05:59}

\begin{knowledgebox}{读图：四条理由其实是一条因果链}
先看经济层：训练主要消耗预算，服务承接用户价值；再看组织层：能做大规模预训练的机构很少，但几乎所有 AI 产品都要运行推理；第三条把 post-training 接回推理；最后一条说明推理横跨 application、scheduler、algebra、network 与 hardware。由此得到的教学结论不是“训练不重要”，而是：模型能力只有在一个跨层系统中才会转化为稳定体验。
\end{knowledgebox}

\begin{warningbox}{不能把 revenue center 理解成“推理一定赚钱”}
推理只有在质量、利用率、价格与可靠性同时成立时才可能创造正向单位经济。低利用率 GPU、过度配置、失败重试、长输出和不可控 agent loop 都会迅速吞掉收入；这里讨论的是价值交付位置，不是财务保证。
\end{warningbox}

\subsection{经验来源与证据边界}

生产经验很有价值，但经验性建议必须标出适用范围。讲者的视角来自 Modal 的 serverless GPU 平台、对数百或数千张 GPU 的客户部署支持、内部 profiling/metrics 工作，以及自己运行推理服务。这解释了为什么本讲特别强调 bursty traffic、cold start、host overhead 与 observability，也提醒读者不要把单个平台的统计量当成所有数据中心的物理定律。

\lecturefigure{slide-004.jpg}{讲者的生产推理经验背景}{官方 deck 第 4 页；课堂 00:06:00--00:07:35}

\begin{warningbox}{来源边界}
本讲义把 slide 上的数字视为讲者在特定生产环境中的经验摘要。尤其是 provider price、GPU failure rate 与 speedup，必须连同 workload、hardware、software version、baseline 和 measurement window 一起解释；缺少这些条件时，只能作为设计提示，不能外推成通用承诺。
\end{warningbox}

\subsection{本章小结}

Inference 的难点来自端到端约束：应用定义正确性和体验，系统把它们变成资源与调度问题，硬件和软件共同决定可行边界。后面的所有技巧都服从同一顺序：\textbf{先定义任务与测量，再选择架构，最后优化实现。}

